\begin{afterword}
\summaryhead

The post opens with a story from a science-fiction show: a captain pries the dematerializing lever off a captured alien ship and expects it to work on his own. Programs that store facts as ``(apple is-a fruit)'' make the same mistake. The word ``apple'' works between people because it pulls a lever on a large shared machinery of seeing, recognizing and handling apples. Written into an AI's knowledge base, it is only the lever. The mistake is tempting because words vary and can be seen, while the machinery is hidden and always present.

The same mistake, the author argues, lies behind proposals to make an AI kind by giving it loving parents. In biology a response that depends on the environment needs more genetic machinery than the response alone. Children absorb culture, and even create grammar, through such inherited responses. An AI has none unless they are designed. Programming a conditional kindness is harder than programming kindness outright, and a paperclip maximizer raised with love stays a paperclip maximizer.

\responsehead

The post's central point is sound and old. A symbol in a program means something to its readers, not to the program. The post credits Drew McDermott's 1976 paper, which mocked the semantic networks of its day. Stevan Harnad later named the same issue the symbol grounding problem, asking how a symbol system's meaning can be more than ``parasitic on the meanings in our heads.''

The biological principle behind the second half is also sound, and the literature on plasticity supports it. A 1998 review of its costs names ``the sensory and regulatory machinery needed for plasticity'' first. The conclusion about AI follows from the principle, as a conditional: an AI without the child's conditional responses will not be changed by being raised as a child.

One example is stated more firmly than the evidence allows. The creole case, Bickerton's bioprogram hypothesis, is given without a hedge, and it is disputed: Siegel argued that the Hawaiian creole grew over two generations, from an expanded pidgin, with grammar from its speakers' other languages. Nicaraguan Sign Language, created by deaf children, supports the post's point better.

``As I've often heard proposed'' points to no one, but the proposal has a long history. Turing suggested educating a ``child machine'' with rewards and punishments in 1950, and described the child brain as ``Rather little mechanism, and lots of blank sheets,'' the view the post calls a chimera. Turing also said that a good child program would have to be found by experiment, which agrees with the post.

The closing claim is the broadest and the least supported. People near AI, the post says, ``usually'' build imitation levers. No case is given, and Robin Hanson objected at the time that this slandered a generation of researchers, citing McDermott's own verdict that most of them did good work. The author replied that ``a very substantial fraction did so, including leaders of the field.''

In short: a sound point about symbols without the machinery behind them, credited to McDermott, carried to AI by a principle from biology that holds up, with a disputed example from creole languages, and ending in a claim about AI researchers that names no case.
\end{afterword}
